% GENERATED FILE. Edit canonical lessons, then run npm run generate:books.
% canonical-source-hash: fnv1a64:aa4675ed673a67f4
% canonical-lessons: GU-C531-jhanjhat, GU-C531-andharyu, GU-C531-ocintu, GU-C531-pherphar, GU-C531-bandobast

\chapter{When Plans Go Wrong}
\label{ch:gu-when-plans-go-wrong}
\hlchaptermodality{531}

\begin{chapteropening}
\textbf{By the end of this chapter:} \emph{I can say what we will do if something unexpected happens.}

Complete the last lesson of chapter 531: I can say what we will do if something unexpected happens.
\end{chapteropening}

\section[\texorpdfstring{\gu{ઝંઝટ} (jhanjhaṭ)}{ઝંઝટ (jhanjhaṭ)}]{\gu{ઝંઝટ} (jhanjhaṭ) — a hassle}

\label{lesson:GU-C531-jhanjhat}



\begin{quote}
Before the new one: say the Gujarati for to push back, then the Gujarati for at long last.
\end{quote}

\subsection*{You'll want to know: \gu{ઝંઝટ}}
\textbf{\gu{ઝંઝટ}} — \emph{jhanjhaṭ} — \textquotedblleft{}a hassle\textquotedblright{}.

\begin{quote}
\gu{જો} \gu{ટ્રેન} \gu{મોડી} \gu{હશે} \gu{તો} \gu{ઘણી} \gu{ઝંઝટ} \gu{થશે}.
\end{quote}

\begin{quote}
\emph{jo ṭren moḍī haśe to ghaṇī jhanjhaṭ thaśe.}
\end{quote}

\begin{quote}
If the train is late, there will be a lot of hassle.
\end{quote}

A feminine noun.

\begin{cousinweb}
\textbf{\gu{ઝંઝટ}} (\emph{jhanjhaṭ}) sounds like what it means, a tangle of small troubles. CDIAL lists Old Gujarati \emph{jhaṁjha}, \textquotedblleft{}a quarrel\textquotedblright{}, and Gujarati \emph{jhā̃jh}, \textquotedblleft{}rage\textquotedblright{}, under Sanskrit \emph{jhañjhā}, \textquotedblleft{}the noise of rain\textquotedblright{} (CDIAL 5326); it does not list this word, so its history is not settled.
\end{cousinweb}

\begin{grammarlens}[title={What you've built}]
One more word for saying what happens if.
\end{grammarlens}

\subsection*{Your turn}
\begin{itemize}
\raggedright
  \item \emph{Say it:} \emph{jhanjhaṭ}
  \item \emph{Say it:} \emph{jhanjhaṭ}, once more
  \item \emph{Say it:} \emph{jo ṭren moḍī haśe to ghaṇī jhanjhaṭ thaśe.} — If the train is late, there will be a lot of hassle.
  \item \emph{Recall:} say the Gujarati for to push back, then the Gujarati for at long last, then say \emph{jhanjhaṭ} again
\end{itemize}

\begin{tcolorbox}[breakable,colback=teal!4,colframe=teal!35!black,title={Before you move on}]
What does \gu{ઝંઝટ} mean? (\textquotedblleft{}A hassle\textquotedblright{}.) Say it once more.
\end{tcolorbox}
\section[\texorpdfstring{\gu{અણધાર્યું} (aṇdhāryũ)}{અણધાર્યું (aṇdhāryũ)}]{\gu{અણધાર્યું} (aṇdhāryũ) — unexpected}

\label{lesson:GU-C531-andharyu}



\begin{quote}
Before the new one: say the Gujarati for at long last, then the Gujarati for a hassle.
\end{quote}

\subsection*{You'll want to know: \gu{અણધાર્યું}}
\textbf{\gu{અણધાર્યું}} — \emph{aṇdhāryũ} — \textquotedblleft{}unexpected\textquotedblright{}.

\begin{quote}
\gu{જો} \gu{કંઈક} \gu{અણધાર્યું} \gu{બનશે} \gu{તો} \gu{મને} \gu{તરત} \gu{ફોન} \gu{કરો}.
\end{quote}

\begin{quote}
\emph{jo kaīk aṇdhāryũ banśe to mane tarat phon karo.}
\end{quote}

\begin{quote}
If something unexpected happens, phone me at once.
\end{quote}

\begin{cousinweb}
\textbf{\gu{અણધાર્યું}} (\emph{aṇdhāryũ}) is built from \emph{aṇ-}, "un-", and \textbf{\gu{ધાર્યું}}, the past of \textbf{\gu{ધારવું}}, \textquotedblleft{}to suppose\textquotedblright{}: not supposed, unexpected. CDIAL files \textbf{\gu{ધારવું}} under Sanskrit \emph{dhāráyati}, \textquotedblleft{}holds\textquotedblright{} (CDIAL 6791).
\end{cousinweb}

\begin{grammarlens}[title={What you've built}]
One more word for saying what happens if.
\end{grammarlens}

\subsection*{Your turn}
\begin{itemize}
\raggedright
  \item \emph{Say it:} \emph{aṇdhāryũ}
  \item \emph{Say it:} \emph{aṇdhāryũ}, once more
  \item \emph{Say it:} \emph{jo kaīk aṇdhāryũ banśe to mane tarat phon karo.} — If something unexpected happens, phone me at once.
  \item \emph{Recall:} say the Gujarati for at long last, then the Gujarati for a hassle, then say \emph{aṇdhāryũ} again
\end{itemize}

\begin{tcolorbox}[breakable,colback=teal!4,colframe=teal!35!black,title={Before you move on}]
What does \gu{અણધાર્યું} mean? (\textquotedblleft{}Unexpected\textquotedblright{}.) Say it once more.
\end{tcolorbox}
\section[\texorpdfstring{\gu{ઓચિંતું} (ocintũ)}{ઓચિંતું (ocintũ)}]{\gu{ઓચિંતું} (ocintũ) — sudden}

\label{lesson:GU-C531-ocintu}



\begin{quote}
Before the new one: say the Gujarati for a hassle, then the Gujarati for unexpected.
\end{quote}

\subsection*{You'll want to know: \gu{ઓચિંતું}}
\textbf{\gu{ઓચિંતું}} — \emph{ocintũ} — \textquotedblleft{}sudden\textquotedblright{}.

\begin{quote}
\gu{જો} \gu{ઓચિંતો} \gu{વરસાદ} \gu{આવશે} \gu{તો} \gu{મેચ} \gu{બંધ} \gu{થઈ} \gu{જશે}.
\end{quote}

\begin{quote}
\emph{jo ocinto varsād āvśe to mec bandh thaī jaśe.}
\end{quote}

\begin{quote}
If a sudden shower comes, the match will be stopped.
\end{quote}

\begin{cousinweb}
\textbf{\gu{ઓચિંતું}} (\emph{ocintũ}) goes back to a reconstructed \emph{*apacitta}, \textquotedblleft{}without thought\textquotedblright{}, and CDIAL lists the Gujarati word, \textquotedblleft{}suddenly\textquotedblright{} (CDIAL 420): what comes without being thought of. It agrees with its noun: \textbf{\gu{ઓચિંતો} \gu{વરસાદ}}.
\end{cousinweb}

\begin{grammarlens}[title={What you've built}]
One more word for saying what happens if.
\end{grammarlens}

\subsection*{Your turn}
\begin{itemize}
\raggedright
  \item \emph{Say it:} \emph{ocintũ}
  \item \emph{Say it:} \emph{ocintũ}, once more
  \item \emph{Say it:} \emph{jo ocinto varsād āvśe to mec bandh thaī jaśe.} — If a sudden shower comes, the match will be stopped.
  \item \emph{Recall:} say the Gujarati for a hassle, then the Gujarati for unexpected, then say \emph{ocintũ} again
\end{itemize}

\begin{tcolorbox}[breakable,colback=teal!4,colframe=teal!35!black,title={Before you move on}]
What does \gu{ઓચિંતું} mean? (\textquotedblleft{}Sudden\textquotedblright{}.) Say it once more.
\end{tcolorbox}
\section[\texorpdfstring{\gu{ફેરફાર} (pherphār)}{ફેરફાર (pherphār)}]{\gu{ફેરફાર} (pherphār) — a change}

\label{lesson:GU-C531-pherphar}



\begin{quote}
Before the new one: say the Gujarati for unexpected, then the Gujarati for sudden.
\end{quote}

\subsection*{You'll want to know: \gu{ફેરફાર}}
\textbf{\gu{ફેરફાર}} — \emph{pherphār} — \textquotedblleft{}a change\textquotedblright{}.

\begin{quote}
\gu{જો} \gu{યોજનામાં} \gu{કોઈ} \gu{ફેરફાર} \gu{થશે} \gu{તો} \gu{હું} \gu{તમને} \gu{કહીશ}.
\end{quote}

\begin{quote}
\emph{jo yojnāmā̃ koī pherphār thaśe to hũ tamne kahīś.}
\end{quote}

\begin{quote}
If there is any change in the plan, I will tell you.
\end{quote}

A masculine noun.

\begin{cousinweb}
\textbf{\gu{ફેરફાર}} (\emph{pherphār}) is built on \emph{pher}, \textquotedblleft{}a turn\textquotedblright{}, which CDIAL files under \emph{*phirati}, \textquotedblleft{}turns\textquotedblright{} (CDIAL 9078), doubled with a rhyming \emph{phār}.
\end{cousinweb}

\begin{grammarlens}[title={What you've built}]
One more word for saying what happens if.
\end{grammarlens}

\subsection*{Your turn}
\begin{itemize}
\raggedright
  \item \emph{Say it:} \emph{pherphār}
  \item \emph{Say it:} \emph{pherphār}, once more
  \item \emph{Say it:} \emph{jo yojnāmā̃ koī pherphār thaśe to hũ tamne kahīś.} — If there is any change in the plan, I will tell you.
  \item \emph{Recall:} say the Gujarati for unexpected, then the Gujarati for sudden, then say \emph{pherphār} again
\end{itemize}

\begin{tcolorbox}[breakable,colback=teal!4,colframe=teal!35!black,title={Before you move on}]
What does \gu{ફેરફાર} mean? (\textquotedblleft{}A change\textquotedblright{}.) Say it once more.
\end{tcolorbox}
\section[\texorpdfstring{\gu{બંદોબસ્ત} (bandobast)}{બંદોબસ્ત (bandobast)}]{\gu{બંદોબસ્ત} (bandobast) — arrangements (made in advance)}

\label{lesson:GU-C531-bandobast}



\begin{quote}
Before the new one: say the Gujarati for sudden, then the Gujarati for a change.
\end{quote}

\subsection*{You'll want to know: \gu{બંદોબસ્ત}}
\textbf{\gu{બંદોબસ્ત}} — \emph{bandobast} — \textquotedblleft{}arrangements (made in advance)\textquotedblright{}.

\begin{quote}
\gu{જો} \gu{મહેમાનો} \gu{રાત્રે} \gu{રોકાશે} \gu{તો} \gu{સૂવાનો} \gu{બંદોબસ્ત} \gu{કરવો} \gu{પડશે}.
\end{quote}

\begin{quote}
\emph{jo mahemāno rātre rokāśe to sūvāno bandobast karvo paḍśe.}
\end{quote}

\begin{quote}
If the guests stay the night, sleeping arrangements will have to be made.
\end{quote}

A masculine noun.

\begin{cousinweb}
\textbf{\gu{બંદોબસ્ત}} (\emph{bandobast}) is usually traced to Persian \emph{band-o-bast}, \textquotedblleft{}binding and tying\textquotedblright{}, and so \textquotedblleft{}arrangements\textquotedblright{}.
\end{cousinweb}

\begin{grammarlens}[title={What you've built}]
One more word for saying what happens if.
\end{grammarlens}

\subsection*{Your turn}
\begin{itemize}
\raggedright
  \item \emph{Say it:} \emph{bandobast}
  \item \emph{Say it:} \emph{bandobast}, once more
  \item \emph{Say it:} \emph{jo mahemāno rātre rokāśe to sūvāno bandobast karvo paḍśe.} — If the guests stay the night, sleeping arrangements will have to be made.
  \item \emph{Recall:} say the Gujarati for sudden, then the Gujarati for a change, then say \emph{bandobast} again
\end{itemize}

\begin{tcolorbox}[breakable,colback=teal!4,colframe=teal!35!black,title={Before you move on}]
What does \gu{બંદોબસ્ત} mean? (\textquotedblleft{}Arrangements (made in advance)\textquotedblright{}.) Say it once more.
\end{tcolorbox}
